\documentclass[11pt]{article}

\usepackage[a4paper,margin=1in]{geometry}
\usepackage[T1]{fontenc}
\usepackage[utf8]{inputenc}
\usepackage{lmodern}
\usepackage{hyperref}
\usepackage{listings}
\usepackage{xcolor}
\setlength{\emergencystretch}{2em}

\lstdefinestyle{embgen}{
  basicstyle=\ttfamily\small,
  breaklines=true,
  columns=fullflexible,
  keepspaces=true,
  frame=single,
  framerule=0.3pt,
  rulecolor=\color{black!35}
}

\title{Embedded Generators: Situated Code Generation for Self-Maintaining Codebases}
\author{Jayanta Majumder \and Sambuddha Majumder}
\date{}

\begin{document}
\maketitle

\begin{abstract}
Modern software systems often contain many layers whose structure is derived
from a smaller set of schemas: database tables, domain models, DTOs, DAOs,
serialization code, user-interface forms, and test fixtures.  When the schema
changes, these derived fragments must change with it.  In practice, the work is
mechanical but distributed, and therefore error-prone.  This paper describes
the embedded generator pattern implemented by \texttt{embgen}, a lightweight
tool that places small generator programs directly inside source-file comments
and rewrites only explicitly marked generated regions.  The approach recovers
some of the practical value of macro systems--code situated next to the
abstraction that produces it--without requiring the host language to be
homoiconic or to provide compile-time procedural macros.
\end{abstract}

\section{Introduction}

The history of programming languages contains a recurring aspiration: programs
should be able to construct other programs when the construction is more
regular, more reliable, or more expressive than hand-written repetition.  Lisp
made this idea unusually direct.  In Paul Graham's account of what distinguished
Lisp, one of the central properties is that ``the whole language always
available'' and that there is ``no real distinction between read-time,
compile-time, and runtime'' \cite{graham-diff}.  He connects this immediately
to macros: ``Running code at compile-time is the basis of macros''
\cite{graham-diff}.  In \emph{On Lisp}, he gives the more operational point:
``A macro generating its expansion has at its disposition the full power of
Lisp'' \cite{graham-onlisp}.

This matters because a macro is not merely a textual abbreviation.  It is a
program that runs in a privileged phase and emits program text, or more
precisely program structure, for a later phase.  The user can create a local
abstraction whose expansion is computed by the same language in which the rest
of the system is written.  In a Lisp-family language, the representation of
programs as ordinary symbolic data gives this technique an unusually natural
form.

Most industrial code, however, is not written in such languages.  C and C++
provide a preprocessor, but its macros are fundamentally token replacement
facilities: a macro name is replaced by its body, parameters are substituted as
preprocessing tokens, and further operators such as stringizing and token
concatenation are available \cite{gcc-cpp-macros,gcc-cpp-args}.  Conditional
compilation is useful, but the facility is not the same as running an arbitrary
host-language program over a rich model and emitting a structured expansion.
Rust is a notable modern exception: its procedural macros run during
compilation and operate over token streams, consuming and producing Rust syntax
\cite{rust-proc-macros}.  But for many languages and many codebases, such
compile-time power is absent, inconvenient, or tied to a particular build
system.

The proposed embedded generators are based on the insight that the practical power of
situated code generation does not require the generated language itself to be
homoiconic.  Nor does it require the compiler to host the generator.  What is
needed is a stable source of truth, a small program that maps that source into a
repeated code pattern, and a set of disciplined locations in the source file where the
generated result is allowed to appear.  XML, JSON, or another schema language
can provide the input representation; comments can host the generator; and a
command-line tool can perform the rewrite before compilation, review, or
deployment.

\section{The Maintenance Problem}

Large applications tend to accrete parallel representations of the same domain
facts.  A field such as \texttt{email} or \texttt{dateOfBirth} may appear in a
database definition, a domain model, a DTO, a DAO, a JSON serializer, a form,
and several test fixtures.  Adding or renaming the field is conceptually simple,
but the edit fans out across many files.  The resulting failures are usually
not intellectually interesting: a mapper omits a field, a fixture drifts from a
schema, a view fails to expose a new property, or a migration no longer matches
the model layer.

Whole-file generators address part of this problem, but they often force a
sharp division between generated and hand-written code.  That division can be
awkward where the repeated fragment is small and embedded inside otherwise
hand-maintained logic.  Developers then face an unattractive choice: either
generate more of the file than they really want to give up, or keep the small
repeated region by hand.

\section{The Embedded Generator Pattern}

\texttt{embgen} uses a narrower contract.  A source file contains a generator
note and a generated region, both marked by comments.  A UUID binds the note to
the region it owns:

\begin{lstlisting}[style=embgen]
//embgen_embedded_generator xml_driven_macro 11111111-1111-1111-1111-111111111111
// @types.xml {/types/type[@name='Person']/fields/field} {
//     emit "    private $javatype $name; // $dbtype\n";
// }
//embgen_generated_start 11111111-1111-1111-1111-111111111111
// generated code appears here
//embgen_generated_end 11111111-1111-1111-1111-111111111111
\end{lstlisting}

The marker syntax is deliberately comment-based.  It can be embedded in Java,
C++, SQL, Ruby, shell scripts, LaTeX, and other formats with single-line
comments.  The host compiler or interpreter treats the generator note as a
comment.  \texttt{embgen}, run as a separate tool, recognizes the marker,
strips the comment prefix from the note, executes the named generator, and
replaces only the text between \texttt{embgen\_generated\_start} and
\texttt{embgen\_generated\_end}.

The UUID is important engineering machinery.  It prevents accidental coupling
to source order and lets multiple generators coexist in one file.  The
surrounding code remains ordinary hand-written source.  The generated region is
small, explicit, and reviewable.

\section{Execution Model}

The Tcl implementation, \texttt{embgen.tcl}, and the Ruby implementation,
\texttt{embgen.rb}, follow the same conceptual pipeline:

\begin{enumerate}
  \item scan each file for embedded generator headers;
  \item collect the note body up to the matching generated-start marker;
  \item dispatch the note to the generator type named in the header;
  \item execute the generator with helper operations such as \texttt{emit};
  \item replace the old generated body with the new emitted text; and
  \item flush any queued secondary files produced by \texttt{emit\_file} or
        \texttt{emit\_to\_file}.
\end{enumerate}

The built-in generator types include schema-driven generators
(\texttt{xml\_driven\_macro} and \texttt{json\_driven\_macro}), procedural
generators (\texttt{tcl\_macro} in the Tcl version and \texttt{ruby\_macro} in
the Ruby version), documentation and diagram generators
(\texttt{dot}, \texttt{plantuml}, \texttt{plantuml\_ascii},
\texttt{latex}, and \texttt{latex\_inline}), an \texttt{echo} generator, and a
\texttt{using\_command\_line} generator.

For XML-driven generation, a note contains one or more directives of the form:

\begin{lstlisting}[style=embgen]
@types.xml {/types/type[@name='Person']/fields/field} {
    emit "private $javatype $name; // $dbtype\n";
}
\end{lstlisting}

The XPath selects nodes.  For each selected node, its attributes are bound as
variables, and the embedded body emits text.  In the Ruby implementation, the
same idea is expressed with Ruby interpolation and REXML node access:

\begin{lstlisting}[style=embgen]
@types.xml {/types/type} {
    table_name = xpathnode.attributes['name']
    emit "CREATE TABLE #{table_name} (\n"
    xpathnode.elements.each_with_index('fields/field') do |field, idx|
        emit ",\n" unless idx.zero?
        emit "    #{field.attributes['name']} #{field.attributes['dbtype']}"
    end
    emit "\n);\n\n"
}
\end{lstlisting}

Relative schema paths are resolved relative to the processed file, walking up
parent directories when needed.  This keeps generator notes relocatable with
their source files while still allowing a shared schema to serve as the project
source of truth.

\section{Examples from the Test Suite}

The repository fixtures exercise the pattern at several scales.  The
\texttt{sample.txt} fixtures demonstrate basic generators: echoing note text,
Graphviz and PlantUML output, XML and JSON field generation, LaTeX image
generation, and command-line output embedding.  The schema fixtures
\texttt{types.xml} and \texttt{types.json} describe domain types such as
\texttt{Person}, \texttt{Address}, \texttt{Order}, \texttt{Product},
\texttt{Invoice}, and \texttt{Tenant}.  From these, the SQL fixtures generate
complete \texttt{CREATE TABLE} statements.

The multi-file SQL fixture demonstrates a more substantial use case.  A single
embedded generator iterates over all schema types and emits one SQL file per
table, using \texttt{emit\_file}.  The generated region in the driver file then
contains a short audit trail such as ``generated file: Person.sql'', while the
actual table definitions are written to separate files.  This shows that the
embedded note can be a local command post for a family of derived artifacts, not
only for the text physically inside the current file.

The Rails view fixture pushes the same idea into a user-interface layer.  For
each type, the generator creates \texttt{show.html.erb} and
\texttt{edit.html.erb} files and derives labels and model variable names from
schema fields.  The value is not that generated views are intrinsically
complete; it is that the repetitive field-level scaffolding has an explicit,
rerunnable derivation.

The permutation fixture demonstrates a non-schema use case.  A procedural macro
uses \texttt{permutations} and \texttt{comma\_separate} to generate exhaustive
test cases for different message orderings.  This is the same core pattern:
replace hand-maintained repetition with a small local program whose output is
checked into the source tree and reviewed as ordinary code.

\section{Relation to Macros}

Embedded generators are macro-like in their purpose but not in their phase
integration.  A Lisp macro is part of the language implementation's expansion
process.  A Rust procedural macro is part of compilation and returns syntax to
the compiler.  An embedded generator is an external rewriting step over marked
regions of ordinary files.

This difference is a limitation and a strength.  It is a limitation because the
host compiler does not type-check the generator itself as part of the target
language, and because incremental rebuild tracking must be supplied by project
convention, CI, or build integration.  It is a strength because the technique is
language-agnostic.  The same schema-driven generator style can maintain Java
fields, SQL DDL, Rails views, documentation snippets, and tests.  The generated
artifact need not even be source code.

The deeper commonality is locality.  In a good macro system, the abstraction
appears at the use site and the expansion is close to the program it explains.
In \texttt{embgen}, the generator note is directly adjacent to the region it
owns.  This preserves discoverability: a maintainer reading the generated code
can see how it is regenerated without searching through a distant build script
or generator project.

\section{Reviewability and Idempotence}

The recommended workflow is to commit both the generator note and the generated
region.  This is sometimes controversial: generated code is often excluded from
version control.  For embedded generators, checking in both parts is a useful
discipline.  The generator is small and local; the output is what the compiler,
database, or renderer actually consumes; and the diff after a schema change
becomes a concrete review artifact.  If a new field is added to
\texttt{types.xml}, the regenerated SQL, model fragments, view fragments, or
fixtures appear in the same review.

The generated region is also bounded.  \texttt{embgen} does not claim ownership
of the whole file.  This reduces the social and technical cost of adoption.  A
team can begin with one painful repeated fragment, give it a UUID, and leave the
rest of the file untouched.

\section{Risks and Boundaries}

The pattern is powerful enough to require discipline.  Embedded code execution
has the same security implications as other local code-generation steps.  A
repository that runs \texttt{using\_command\_line}, diagram tools, LaTeX, Tcl, or
Ruby bodies is executing code, not merely formatting text.  Projects should run
generators in trusted contexts, review generator notes carefully, and integrate
the tool into CI only with the same care given to build scripts.

Generator notes should also stay small.  If a generator becomes a large program,
it may deserve to be promoted into a named library or external tool.  The
embedded form is best when it captures a local, regular, reviewable mapping
from schema data or simple procedural logic into a nearby code region.

\section{Conclusion}

\texttt{embgen} generalizes a practical lesson from macro systems: many of the
best abstractions are programs that write the repetitive parts of other
programs.  Lisp offers this through a language whose code representation and
phase structure make macros exceptionally expressive.  Rust offers another
modern route through procedural macros.  Embedded generators offer a third,
more portable route: place the generator in comments, place the schema in XML,
JSON, or another explicit representation, and let a small tool refresh only the
marked generated regions.

The result is not a replacement for true language-integrated macros.  It is a
pragmatic technique for ordinary polyglot codebases, where schema drift and
boilerplate maintenance are everyday costs.  By making the derivation local,
explicit, rerunnable, and reviewable, embedded generators allow codebases to
retain hand-written structure while making their mechanical fragments
self-maintaining.

\begin{thebibliography}{9}

\bibitem{graham-diff}
Paul Graham.
\newblock \emph{What Made Lisp Different}.
\newblock December 2001, revised May 2002.
\newblock \url{https://www.paulgraham.com/diff.html}

\bibitem{graham-nerds}
Paul Graham.
\newblock \emph{Revenge of the Nerds}.
\newblock \url{https://www.paulgraham.com/icad.html}

\bibitem{graham-onlisp}
Paul Graham.
\newblock \emph{On Lisp: Advanced Techniques for Common Lisp}.
\newblock 1993.
\newblock \url{https://www.paulgraham.com/onlisp.html}

\bibitem{gcc-cpp-macros}
Free Software Foundation.
\newblock \emph{The C Preprocessor: Macros}.
\newblock GCC documentation.
\newblock \url{https://gcc.gnu.org/onlinedocs/cpp/Macros.html}

\bibitem{gcc-cpp-args}
Free Software Foundation.
\newblock \emph{The C Preprocessor: Macro Arguments}.
\newblock GCC documentation.
\newblock \url{https://gcc.gnu.org/onlinedocs/cpp/Macro-Arguments.html}

\bibitem{rust-proc-macros}
The Rust Project Developers.
\newblock \emph{The Rust Reference: Procedural Macros}.
\newblock \url{https://doc.rust-lang.org/reference/procedural-macros.html}

\end{thebibliography}

\end{document}
